\section{Context Loop Wars：入口、工具与数据回流}

当 coding agent、办公 agent 或行业 agent 进入真实工作流，掌握用户入口的公司可以观察任务分布，拥有模型的公司可以控制能力与价格，拥有工具和环境的公司则决定行动边界。三种位置不总由同一家公司占据，于是合作关系会迅速转为竞争。本节把讲者选取的 Windsurf、OpenAI、Anthropic 与 Mistral 案例当成系统结构分析，而不是商业八卦。

\subsection{一次收购争夺暴露了什么}

讲者用 2025 年 coding-agent 市场的变化说明 context loop 的价值。Windsurf 位于开发者工作流入口，模型供应商提供底层能力，用户代码与执行反馈构成高价值环境。只要其中一方担心另一方把交互数据、分发或产品关系内化，原本的 API 合作就会出现访问限制、收购谈判和替代模型部署。

\begin{figure}[H]
\centering
\includegraphics[width=0.94\textwidth]{images/00-27-34-context-loop-wars.jpg}
\caption{Context Loop Wars 的第一阶段：围绕 Windsurf 工作流入口与模型访问产生冲突。}
\figsource{Stanford Online 官方视频 00:27:34，`images/00-27-34-context-loop-wars.jpg`。}
\end{figure}

\begin{knowledgebox}{读图：不要只盯着收购价格}
更重要的是四条流：用户请求从哪里进入，代码和工具在何处执行，哪家模型提供推理，结果与修改回到谁的数据系统。入口公司掌握需求和分发，模型公司掌握能力与供给，云或工具公司掌握执行基础设施。交易只是重新安排这四条流的控制权。理解结构后，即使具体公司变化，分析方法仍然有效。
\end{knowledgebox}

\begin{figure}[H]
\centering
\includegraphics[width=0.94\textwidth]{images/00-29-04-context-loop-wars-cont.jpg}
\caption{Context Loop Wars 的延伸：模型、云、开发工具与人才在同一闭环中重新组合。}
\figsource{Stanford Online 官方视频 00:29:04，`images/00-29-04-context-loop-wars-cont.jpg`。}
\end{figure}

第二张图把单次交易扩展为生态关系。读图时应比较每个参与者缺少什么：模型实验室可能缺稳定分发与高质量环境，开发工具可能缺自有模型和长期推理成本控制，云平台可能希望把算力需求锁入自身基础设施，人才与研究团队则会随着资本和数据位置迁移。所谓 context war，实质是争夺可持续改进闭环，而不是争夺一批静态聊天记录。

\begin{warningbox}{课堂提示：案例说明机制，不保证时间线永远有效}
并购状态、API 政策和产品能力都会变化。讲义保留的是讲者在该次课程中用来解释 context 价值的历史案例。阅读者应验证最新事实，再用“入口—环境—模型—反馈—分发”框架判断新局面，不能把课件中的公司连线当作永久产业地图。
\end{warningbox}

\subsection{Mistral：关键 context 也可以是主权环境}

Context 不只属于某个 SaaS 产品。语言、法律、公共部门采购、国家安全要求、数据驻留和本地产业生态也会形成难以跨地区复制的环境。讲者把 Mistral 的创建动机放在欧洲主权与开放权重需求中解释：当关键模型能力完全依赖外部供应商时，本地机构很难控制数据、成本、可审计性与长期议价权。

\begin{figure}[H]
\centering
\includegraphics[width=0.94\textwidth]{images/00-30-00-mistral-critical-contexts.jpg}
\caption{Mistral 案例：欧洲语言、法规、公共机构和产业需求构成 critical contexts/environments。}
\figsource{Stanford Online 官方视频 00:30:00，`images/00-30-00-mistral-critical-contexts.jpg`。}
\end{figure}

\begin{knowledgebox}{读图：sovereign AI 的系统含义}
Sovereign AI 不只是“训练一个国产模型”。它至少涉及本地语言与文化数据、可接受的模型许可、数据驻留、供应链控制、公共部门部署接口、人才与算力，以及在危机时保持服务的能力。开放权重可以提高可检查性和部署选择，却不会自动解决训练成本、安全责任或生态碎片化。主权诉求最终仍要落到可运行的基础设施与制度安排。
\end{knowledgebox}

\begin{importantbox}{Context 的三种稀缺性}
产品 context 稀缺，因为真实工作流和用户信任难以复制；组织 context 稀缺，因为私有知识、权限和历史决策不能公开抓取；主权 context 稀缺，因为语言、法律与公共利益具有地域性。模型能力越接近，谁能安全进入这些环境，谁就更可能获得下一轮高价值反馈。
\end{importantbox}

\subsection{本章小结}

Context loop wars 把抽象的“数据优势”具体化为入口、工具、执行环境、模型供给和反馈回流的控制权。Windsurf 案例展示产品闭环，Mistral 案例展示主权与公共环境。下一节继续追问：即便拥有环境和反馈，RL 是否真的能够持续扩展，还是会遇到可测量与不可测量的边界。

